\section{R network workflows: the exposure, assignment and population must travel together}
\label{sec:exec25-r-network}

A bank choosing which branches receive a retention intervention faces several different network questions. It may ask for an exposure-specific response under a randomized allocation, an effect among customers whose treatment changes with an encouragement, a test of a sharp exposure null, or the value of a policy selected from a constrained class. These are different targets. A common software interface must preserve their distinctions instead of treating every object called a spillover effect as an interchangeable number. The network itself is part of the observation contract: borrower identifiers, direction, weights, missing relationships and the time at which an edge became available determine which exposure is being measured.

The present comparison extends the earlier inferference review with eight R implementations at authentic pre-cutoff revisions. The selected versions are interference 0.1.0, latenetwork 1.0.1, BicliqueRT 1.0.0, NetworkTargeting 0.0.0.9000, NetworkCausalTree 1.0.0, MisspecifiedInterference 1.2, autognet 0.0.3 and clusteredinterference 1.0.1. Their inspected revisions range from 2019 through June 2026 and precede the client cutoff of 7 October 2026. Repository dates and declarations identify source bytes; none of these selected packages or their tests was executed in this comparison. The findings below concern inspected expressions, composed call paths and project-authored exact witnesses. They are not installed-package benchmarks or imported proofs from papers named in source comments.~\citep{exec24_interference_inferference,exec25_rnet_exposure,exec25_rnet_late,exec25_rnet_biclique,exec25_rnet_target,exec25_rnet_tree,exec25_rnet_nmr,exec25_rnet_auto,exec25_rnet_cluster}

\subsection{Design-based exposures: a label is a mathematical contract}

Let $Z=(Z_1,\ldots,Z_n)$ denote the assignment vector, $A$ the declared adjacency matrix, and $e_i(Z,A)$ a specified exposure map for customer $i$. A mean under exposure $c$ requires an identifying relationship between the potential response and that map, as well as support under the assignment law. The notation $\pi_i(c)=\Pr\{e_i(Z,A)=c\}$ concerns that law; it is not an ordinary propensity inferred merely from the presence of a network field. Joint probabilities $\pi_{ij}(c,d)$ also matter for uncertainty. A command must retain the design or simulator, the map, and the precise population along with the resulting estimates.

The interference package constructs first-degree categories from own assignment and a positive peer-exposure indicator. Its selected expression multiplies the assignment row vector by the adjacency matrix. Thus an asymmetric adjacency convention must be checked against the bank's arrow convention. Its two-hop option uses a thresholded matrix square with a zeroed diagonal. A two-step path is not necessarily a neighbor at shortest-path distance exactly two: triangles can create overlap with direct neighbors. A translation layer must describe that convention rather than silently replace it with a different graph operator.~\citep{exec25_rnet_exposure}

One inspected full-neighborhood helper is particularly useful as a qualification target. Its documentation names two conditions: the focal unit and all neighbors treated, or the focal unit and all neighbors controlled. The selected implementation instead defines the second indicator by exposure count less than neighborhood size. In a two-node connected graph with assignment $(1,0)$, neither closed neighborhood is all controlled, yet both are not all treated. The authored witness therefore produces two true selected control indicators and two false literal all-control indicators. The distinction is substantive: a comparison with a mixture of partially treated neighborhoods answers a different question from comparison with universal local control. Admission should wait for an explicit choice of contract and a version-specific native test.~\citep{exec25_rnet_exposure,exec25_rnet_choice}

The probability builder also makes an explicit approximation. With $R$ simulated assignments and exposure-indicator matrix $I_c$, the same-exposure matrix is
\[
 \widehat P_{cc}=\frac{I_cI_c'+I_n}{R+1},
 \qquad
 \widehat P_{cd}=\frac{I_cI_d'}{R}\quad(c\ne d).
\]
The additive diagonal correction is not the exact enumerated design probability. If four assignments exhaust a two-person Bernoulli design and each of four categories occurs once for a focal person, the four corrected marginals are $2/5$ and sum to $8/5$, whereas the exact probabilities are $1/4$. This does not by itself invalidate every use of a deliberate Monte Carlo correction. It requires the adapter to distinguish corrected finite-simulation quantities from an exact probability distribution and to assess their effect on estimates and uncertainty. A positive correction for an unseen event does not establish that the scientific design supports that event.~\citep{exec25_rnet_exposure,exec25_rnet_choice}

The selected permutation generator rounds $np$ to a fixed treated count, samples assignments and rejects duplicates unless repetition is allowed. That is a fixed-count design, not independent Bernoulli assignment. The estimator then distinguishes a Horvitz--Thompson total from a H\'ajek normalized total, and marks an exposure with no observed contributors as unavailable. The proposed reuse is a restricted design-based adapter with explicit mapping and probability modes. Direct adoption of a generic exposure label without these choices is rejected. Complete covariance helpers, random-state behavior and dependencies remain qualification work.~\citep{exec25_rnet_exposure,exec25_rnet_choice}

\subsection{Instrumented network responses: support and the variance law remain separate}

The latenetwork direct path forms empirical joint strata of an instrument $Z_i$ and an instrumental exposure $T_i$. It compares conditional outcome and treatment means across instrument values, then returns their ratio. Schematically the ratio is
\[
 \widehat\theta(t)=
 \frac{\widehat E[Y\mid Z=1,T=t]-\widehat E[Y\mid Z=0,T=t]}
      {\widehat E[D\mid Z=1,T=t]-\widehat E[D\mid Z=0,T=t]}.
\]
For a bank encouragement, assignment validity, exclusion, treatment-response restrictions and the interpretation of the exposure stratum are substantive requirements. An empty empirical stratum or near-zero denominator is a separate numerical and inferential problem. Neither is resolved by a function returning a local-effect label.~\citep{exec25_rnet_late}

The selected validator rejects missing inputs, nonbinary assignments, asymmetric adjacency and nonzero diagonal entries. In particular, its symmetry test counts unequal transpose entries; it does not sum signed differences that would cancel. The direct routine constructs undirected graph distances on the original population before restricting to the requested subpopulation. Paths through customers outside that subpopulation can therefore remain in the distance matrix. This differs from computing distances on the induced subgraph and must be carried into the uncertainty contract.~\citep{exec25_rnet_late}

The inspected HAC branch uses $K_{ij}=\mathbf 1\{d(i,j)\le b\}$ in a score quadratic, followed by a square root. A distance cutoff matrix need not be positive semidefinite. For the three-node path at bandwidth one and a centered illustrative score,
\[
 K=\begin{pmatrix}1&1&0\\1&1&1\\0&1&1\end{pmatrix},
 \qquad v=(1,-2,1)',\qquad \frac{v'Kv}{9}=-\frac{2}{9}.
\]
This exact matrix witness shows a property of the selected quadratic, not a claim that a particular native fitted sample produces this score. The wild-bootstrap path uses a different neighborhood-overlap covariance and a separate small-eigenvalue rule. These are not interchangeable uncertainty services. Qualification must fix graph distances, bandwidth, score law, weak-instrument handling and all nonfinite results before comparing interval coverage. The proposal retains latenetwork as a bounded reference and adapter candidate, with no universal network-IV or dependence guarantee.~\citep{exec25_rnet_late,exec25_rnet_choice}

\subsection{Conditional randomization tests need typed failure and tie semantics}

BicliqueRT accepts a design function, an exposure function and a null-equivalence relation. The inspected preparation places the observed assignment in the first column, simulates other assignments from the supplied design, and builds a null-exposure graph. The selected test computes a statistic on the focal units and assignments of the chosen biclique. The native bimax route has its own compiled call and bounded result capacity. These are useful building blocks for a declared conditional test; the package title does not establish validity for an arbitrary assignment sampler or an incorrectly specified null.~\citep{exec25_rnet_biclique}

The return contract needs special attention. Missing statistics are returned with a numerical p-value field equal to three; a degenerate randomization distribution can return two. These values represent exceptional states, not probabilities. A CLI or MCP wrapper must translate them into explicit unavailable-test statuses while retaining the original diagnostic. A downstream agent that merely compares the field with $0.05$ could otherwise describe a failed test as evidence of no effect.~\citep{exec25_rnet_biclique}

The selected randomized-tail helper counts approximate ties with $|t_r-t_0|<10^{-14}$, but separately counts all strict exceedances $t_r>t_0$. Those two sets can overlap. Its expression is
\[
 p_U=\frac{\#\{r:t_r>t_0\}+U(1+\#\{r:|t_r-t_0|<\varepsilon\})}{R+1}.
\]
Take $t_0=0$, simulated values $(10^{-15},10^{-15},1)$, $\varepsilon=10^{-14}$ and the stipulated draw $U=3/4$. There are three strict exceedances and two approximate ties, giving $p_U=21/16$. The two-sided expression $2\min(p_U,1-p_U)$ then gives $-5/8$. The nondegeneracy condition is satisfied in this authored example. This is a counterexample for the selected arithmetic contract, not a native package run. A proposed repair must partition ties and strict tails consistently and then prove or test the intended randomization law; clipping a number into $[0,1]$ is not a validity argument. Random-state identity, tolerance scaling and refusal behavior belong in the qualification suite.~\citep{exec25_rnet_biclique,exec25_rnet_choice}

\subsection{Policy targeting: numerical optimization does not restore missing graph truth}

NetworkTargeting proposes welfare-maximizing allocations under network features and treatment constraints. Its selected wrapper refuses missing treatment or primary features, but handles other missing inputs by changing the focal population and imputing values. Missing contextual covariates receive column means; missing entries of selected network matrices receive zero; affected rows become nonfocal. A zero is an asserted absence of a relationship, whereas an unknown edge is an observation problem. The resulting policy objective therefore needs its own explicit population and missingness interpretation.~\citep{exec25_rnet_target}

The inspected MILP wrapper forwards the first constraint-tolerance parameter to both tolerance arguments. The solver helper has explicit feasibility settings and a time limit when it constructs its own parameters, then returns the Gurobi result. The selected caller consumes the solution vector and objective in subsequent checks before an explicit termination-status guard in that path. An admissible policy response must distinguish a feasible incumbent, an optimal solution within tolerance, an infeasible request, a time-limited result and missing solution fields. It must also preserve focal population, budget and out-of-sample policy value. A finite assignment vector alone certifies none of those properties.~\citep{exec25_rnet_target}

Source rights also need a concrete resolution. A root GPL version-three text is present, while DESCRIPTION still contains a license placeholder; the selected dependencies include Gurobi. Those observations are not a complete distribution or entitlement decision. The proposal keeps this candidate as a policy-optimization reference pending coherent rights, solver and scientific contracts. Direct embedding or porting is deferred. The known route is to fix the target, trace the complete optimization and score paths, test explicit termination states, and compare a frozen public allocation problem against an exact small reference.~\citep{exec25_rnet_target,exec25_rnet_choice}

\subsection{Heterogeneous network effects: cluster splitting does not justify changing exposure}

NetworkCausalTree builds a discovery tree from sampled clusters and estimates effects on the remaining clusters. Its estimation path constructs the induced adjacency matrix of the retained units and recomputes neighborhood exposure there. When clusters are disconnected components, this can preserve the relevant relationships. If treatment across a discarded cross-cluster tie affects a retained unit, the induced graph changes the exposure definition. Honest sample splitting and scientifically correct exposure construction are separate requirements. The bank must supply a defensible cluster boundary or retain the appropriate background exposure.~\citep{exec25_rnet_tree}

The public argument describes assignment probability as a unit-specific vector. The inspected marginal helper nevertheless computes the chance of at least one treated neighbor using the focal unit's probability raised to its degree. For a focal probability $1/5$ and one neighbor probability $4/5$, independent assignments give
\[
 \Pr(W_i=0,G_i=1)=(1-1/5)(4/5)=16/25,
\]
while the selected expression uses $(1-1/5)(1/5)=4/25$. Equal assignment probabilities avoid this particular discrepancy. The adapter must therefore qualify heterogeneous assignment probabilities explicitly rather than infer support from a vector-shaped argument.~\citep{exec25_rnet_tree,exec25_rnet_choice}

There is a separate adjacent-pair issue to test. In a two-node connected graph with independent probability-one-half assignments, both units having own treatment one and a treated neighbor occurs exactly when both are treated, with probability $1/4$. The selected joint helper expression evaluates to $1/16$ when the shared-neighbor count is zero. The inspected variance path, however, calls its pair correction only when that count is positive: adjacent nodes with no common third neighbor are a separate dependency case. The authored helper witness and the actual pair-selection branch must both be tested; they should not be conflated into an assertion about a completed native interval. The proposal defers native admission pending exposure, marginal, joint-probability, pair-selection and tree-selection qualification.~\citep{exec25_rnet_tree,exec25_rnet_choice}

\subsection{Candidate graphs and auto-g-computation answer different uncertainty questions}

MisspecifiedInterference takes a list of candidate graphs and forms an exposure indicator that requires the same condition under every graph in that list. Its simulated probability matrices use an additive diagonal correction, and its mean routines expose both fixed-population and estimated-denominator forms. This construction makes the candidate list scientifically consequential. If one candidate leaves the focal unit isolated, an indirect-exposure condition can be impossible under that candidate. Requiring agreement with a one-edge alternative can then have probability zero. A diagonal correction of $1/(R+1)$ after zero observed occurrences cannot make that target identifiable. At $R=9$ it is numerically $1/10$, while structural support remains zero.~\citep{exec25_rnet_nmr,exec25_rnet_choice}

The route forward is to separate known design probabilities, simulation error, candidate-graph restrictions and substantive guarantees that the true exposure model is represented. An empty common event requires refusal, a changed target or justified information; it is not a numerical smoothing problem. The selected DESCRIPTION contains an unevaluated license-helper placeholder and the inspected complete repository tree has no license original. This candidate is retained for reference study, with copying or distribution deferred until rights and native behavior are resolved. No general repair for unknown bank relationships is inferred.~\citep{exec25_rnet_nmr}

autognet instead evaluates network response models through parameter draws and nested Gibbs simulations. The inspected R method selects parameter iterations by burn-in and thinning, converts adjacency identities for native calls, and constructs average, direct and spillover contrasts from distinct simulated targets. The native individual-intervention path runs a separate chain for each selected person. Its averaging mode retains sampled binary outcomes after burn-in; the nonaveraging mode returns the final conditional response probability. These are different Monte Carlo outputs even when they concern the same model.~\citep{exec25_rnet_auto}

The number of parameter draws, number of within-network simulation sweeps, initial states, retained iterations and individual intervention population all belong in the result record. More sweeps may reduce simulation error conditional on a model, but do not validate that model or certify chain mixing. The inspected source provides a bounded integration candidate after index, allocation, convergence, Monte Carlo error and dependency tests. An exact small-network enumeration can serve as the first reference before larger public comparisons.~\citep{exec25_rnet_auto,exec25_rnet_choice}

\subsection{Clustered policies: an available-group mean is not an unchanged population}

clusteredinterference supplies a group-based policy workflow with estimated nuisance parameters, treatment-vector evaluation and a sandwich variance. Its selected allocation helper samples distinct vectors with a prescribed treated count and carries a multiplicity factor into the probability integrand. Asking for all vectors sets the target number to the combinatorial support size, but still collects them through random proposals and duplicate rejection. For six equiprobable vectors the expected number of proposals to collect all six is $6\sum_{j=1}^6 1/j=147/10$. There is no deterministic upper bound on draws from that loop. A complete public test should report proposal count and incomplete work, not only the requested number of distinct vectors.~\citep{exec25_rnet_cluster,exec25_rnet_choice}

The quadrature helper integrates a single random effect over five standard deviations on either side of zero. A caught integration error or a missing integral returns numeric zero. Numerical failure and a scientifically zero assignment probability are different states and must remain distinct in the adapter. The selected target routine takes the mean of group contributions with missing values removed. For contributions $(2,\text{missing},8)$, the available-group mean is five; summing available contributions over the original three-group denominator gives $10/3$. Neither calculation supplies the unknown contribution. The distinction matters particularly because earlier inferference paths retain a different denominator convention. Group omissions must be visible before any cross-package comparison.~\citep{exec25_rnet_cluster,exec24_interference_inferference,exec25_rnet_choice}

The sandwich routine solves an assembled bread matrix and sums group-score outer products. Its inputs, derivative choices, singularity behavior and relation to the actual number of usable groups must be reviewed together. The proposal retains a conditional group-policy reference, with numerical-failure typing and population accounting as explicit admission requirements. Cluster labels cannot establish absence of cross-group interference or shared calendar shocks.~\citep{exec25_rnet_cluster}

\subsection{A public qualification portfolio with independent candidate decisions}

The first public fixtures should enumerate tiny assignment spaces and verify exposure categories, support, marginals and joint probabilities exactly. They should include asymmetric graphs, a triangle, adjacent nodes without common neighbors, heterogeneous assignment probabilities, isolated nodes and cross-cluster ties. The second group should test numerical and output contracts: empty exposures, weak instrument denominators, indefinite score kernels, near ties, failed randomization statistics, missing group contributions, quadrature errors, time-limited optimization and nested simulation retention. Every request must retain the scientific target and a typed failure state across library, CLI and MCP.~\citep{exec25_rnet_choice}

The third group concerns evidence beyond arithmetic. Match the assignment and sampling law, compare uncertainty over fresh public or generated studies, retain complete failure denominators, and measure total adapter, compute and review costs. An illustrative USD 30,000 program costs USD 30 per 1,000 attempts, USD 50 per 600 completions and USD 75 per 400 accepted outputs. These exact arithmetic comparisons are hypothetical and do not establish bank returns. A useful public result can justify later bank experiments; it cannot supply missing assignment truth or graph ownership.~\citep{exec25_rnet_choice}

The eight candidates remain separate reuse decisions. Bounded searches contain irrelevant interference applications, discovery-only repositories and uninspected tails. Other observational, simulation, spatial and experiment-design families remain open; the selected originals do not establish exhaustive R coverage. Complete delegates, native artifacts, dependency rights and scientific restrictions remain explicit. A locally unresolved license or source path does not prevent another candidate's qualification from advancing. The proposal therefore records an inconclusive coverage outcome while specifying executable next tests, without claiming complete network qualification.~\citep{exec25_rnet_choice}
